% sections/11_scope.tex -- Section 11, Scope and Sensitivity (W4; 11.2 drafted by W2 and \input below).
% Plan: BLUEPRINT 2.11, 5, 6; Overrides O1, O5, O8 (3.5-page budget). Register: BLUEPRINT 9. Every number carries a "% src:" comment.
% Abbreviations: T: = paper/tables/; F: = paper/figures/abdou/; 5:NN, 11:NN, 12:NN = frot/sub2025/{5_quantitative_model,11_app_data,12_app_quant}.tex;
% v6: = frot/current/; D0N = dossier/0N_*.md.
\section{Scope and Sensitivity}\label{sec:scope}

The reform tables of \autoref{sec:elimination} to \autoref{sec:menu} are computed in the model of \autoref{sec:model} at the estimates of \autoref{tab:est}. This section states the restrictions under which they hold and, for each, the direction in which relaxing it would move the results, in the order of \autoref{subsec:scope}.

\subsection{Statutory ages as reference points}\label{subsec:reference}

In the model the statutory ages act only through the budget constraint and the program rules. Age 62 opens eligibility for retirement benefits and makes pension income liquid. The normal retirement age is the kink of the claiming-age schedule in \eqref{eq:adjust}, the age at which the earnings-test rate and threshold change in \eqref{eq:ret}, and the age at which eligibility for disability insurance ends in \eqref{eq:di}. Age 65 brings Medicare coverage into \eqref{eq:oop}, and age 70 ends the earnings test and makes claiming mandatory. % src: T:parameters_ss.tex:46-52; 12:69-71 (pension liquidity at 62); 12:49 (DI age condition); 5:34-40 (Medicare from 65); 12:97 (claiming by 70)
There are no default, norm, or reference-point effects: an individual who claims at 65 does so because the schedule, the earnings test, and the health and wage profiles make it worthwhile, not because 65 is the normal retirement age. % src: T:parameters_ss.tex:47 (NRA 65)

The evidence indicates that such effects exist. \citet{seibold2021reference} documents, from discontinuities in German pension rules, bunching of retirement at statutory ages that financial incentives alone do not account for. \citet{mastrobuoni2009labor}, \citet{behagel2012}, and \citet{deshpande2024sticky} study the US cohorts whose normal retirement age was raised by the 1983 amendments and find that the spike in claiming and retirement at the normal retirement age moved with it; \citet{lalive2023social} find a similar response to a statutory-age increase in Switzerland.
% src: 1:44 (1983 amendments; footnote of the May 2025 introduction)
\citet{liebman2015would} find that information about the program's rules raised the labor force participation of older workers but did not shift their claiming, and \citet{brown2016framing} and \citet{shoven2017} find that the framing of claiming information and the reasons individuals give for their claiming ages go beyond the actuarial gain from delay.

The paper takes one stance on what this implies. The predictions of \autoref{sec:claiming} for a higher normal retirement age are the incentive component of the response. A reference effect at the new normal retirement age would move claiming later for individuals who, in the model, keep claiming at 62--64 after the reform. It would raise the government's saving, since delayed claims are paid for fewer years at a factor closer to 100; it would lower the benefit cut borne by those who move; and it would add a cost for those who delay for reasons the model does not price. For a steeper actuarial adjustment the effect works the other way. The normal retirement age stays at 65, so individuals anchored to it respond less to the steeper reward and the steeper penalty than the model predicts, and the participation gains at ages 65--69 and the benefit savings in the AD and AL rows of \autoref{tab:claiming} would be overstated. Reference effects therefore move the contrast of \autoref{sec:claiming} toward the normal-retirement-age reform: a higher statutory age moves a focal point, whereas a steeper adjustment changes the schedule without moving it. The welfare sign of the contrast does not change, because the benefit cut under a higher normal retirement age, every factor of \autoref{tab:schedules} lower, does not depend on behavior. % src: T:DRC_visualization_4.tex:38-40
The size of either correction depends on the distribution of claiming ages in the model and in the data; claiming ages are not among the estimation targets and are not reported, so the discussion is qualitative.
% TODO-AUTHORS: Open question 3 -- if the baseline simulation panels are post-processed into a claiming-age distribution (model, and data from the PSID), replace the last sentence with the comparison.

% sections/11_2_factor_prices_draft.tex -- draft of subsection 11.2 (W2b) for W4 to \input into sections/11_scope.tex.
% Plan: BLUEPRINT 2.11 (11.2), 5.1-5.2; Decision on numbers: none reported here (no digits in the text; signs only).
% The algebra and the labeled equations eq:ge-prices, eq:ge-welfare, eq:ge-fiscal, eq:ge-mvpf and the sign table
% tab:gesigns are in appendix/D_derivations.tex (subsection app:factorprices); this text refers to them with \eqref/\autoref.
% W4: drop the \subsection line below if 11_scope.tex supplies its own heading; the label subsec:factorprices is referenced nowhere else yet.
\subsection{Factor prices: partial equilibrium and a reduced-form adjustment}\label{subsec:factorprices}

Every column of the reform tables is computed at the status-quo prices: the gross return $r_f$ and the age--health wage profiles $W(h,t)$ of \autoref{sec:model} do not respond to the cohort's choices. \autoref{subsec:scope} states the assumption as that of a small open economy, in which fixed $r$ and fixed wages are one assumption rather than two (Appendix~\ref{app:factorprices}). The reforms nevertheless change what the cohort supplies to factor markets: the participation columns of Tables~\ref{tab:claiming}--\ref{tab:other} report changes in labor supply by age group, and the saving subsidies change asset holdings (\autoref{subsec:saving}). In a closed economy, or if the reform applied to every cohort of a stationary economy, these changes would move the capital--labor ratio and with it the wage and the return. % src: T:welfare_savings_mvpf_v3.tex:120-123 (Indir. < 0 in every K row); 7:423 (taxes on asset income)

Appendix~\ref{app:factorprices} derives a reduced-form adjustment that requires no new solution of the model. With a constant-returns technology and competitive factor markets, the proportional change in the wage is the capital share divided by the elasticity of substitution, times the change in the capital--labor ratio, and the proportional change in the return gross of depreciation has the opposite sign and the labor share in place of the capital share, as in \eqref{eq:ge-prices}. At fixed behavior, the price changes alter each individual's earnings, net asset income, and, through the earnings record, retirement benefits. Valued at the relative marginal utilities of \autoref{sec:suffstats}, they add a term to the dollar willingness to pay, \eqref{eq:ge-welfare}, and, through the income tax, the payroll tax, and the benefit base, a term to the change in net spending, \eqref{eq:ge-fiscal}; the adjusted MVPF is the ratio of the adjusted numerator to the adjusted denominator, \eqref{eq:ge-mvpf}. The two terms partition the cohort's change in factor income between its members and the government.

The direction of the adjustment follows from the participation columns and is collected in \autoref{tab:gesigns}. Reforms that raise labor supply---the income-tax reductions at every age, and the claiming-schedule reforms after the normal retirement age---lower the wage and raise the return in a closed economy: workers lose part of the gain the tables record, asset holders gain, part of the indirect tax saving is returned, and the benefit base falls. Reforms that lower old-age labor supply---the larger benefit-formula increases and the Medicare expansions---raise the wage and lower the return, with the opposite incidence. The saving subsidies raise the capital stock and lower old-age participation, so they raise the wage and lower the return: the higher return they pay is partly undone, while labor-tax revenue and the benefit base rise. For the claiming-schedule reforms the sign of the aggregate labor-supply change depends on the weights, in efficiency units of labor (Appendix~\ref{app:factorprices}), of age groups whose participation moves in opposite directions. % src: signs of the participation columns in T:welfare_age_dependence_mvpfv5.tex:108-126, T:welfare_progressivity_mvpfv7.tex:115-158, T:welfare_medicare_mvpfv8.tex:113,116, T:welfare_tax_mvpf_v8.tex:114-118, T:welfare_savings_mvpf_v3.tex:120-123

No adjusted MVPF is reported. The participation changes in the tables are at most a few percentage points within age groups that are a minority of the cohort's labor in efficiency units, and the adjustment scales them by a factor share and divides by the substitution elasticity; the paragraph states the direction of the adjustment and not its size. A number would require the capital share, the elasticity of substitution, and the depreciation rate; the reform-induced changes in the cohort's labor in efficiency units and in its assets, computable from the simulated panels behind the tables; the age profiles of the labor-tax, asset-tax, and benefit bases in \eqref{eq:ge-fiscal}; and the elasticity of benefits to the wage profile. The companion paper lists factor-price feedback among the questions it leaves open \citep{ndiaye2026flexible}; the fixed-price environment used here is that of \citet{farhi2013insurance} and \citet{golosov2016redistribution}. % src: v6:9_conclusion_resubmission_v6.tex:10


\subsection{Couples, bequests, health information, time separability, income risk, and reentry}\label{subsec:restrictions}

\paragraph{Couples.}
The decision-maker is an individual, and the estimation targets are profiles for household heads and spouses in the general population (\autoref{tab:psid}). % src: T:sample_psid.tex:10
Spousal and survivor benefits, which are computed from the worker's own benefit, and the coordination of claiming and labor supply within couples are outside the model; \citet{Borella2023} estimate a model of couples in which marriage-related Social Security rules shape the labor supply of both spouses. In such a model a reform of the claiming-age schedule also changes the base of the spousal and survivor benefits, and the claiming age of each spouse responds to the other's; neither margin is in the reform tables.

\paragraph{Bequests.}
Bequests are accidental and assets at the terminal age are zero (\autoref{subsec:budget}). Mean assets are matched through the 40s and fall below the data from the early 50s; median assets are above the data from the 30s through the early 50s and below it from the mid-50s; at age 60 the model has 1.642 against 3.218 for the mean, in hundreds of thousands of 2016 dollars (\autoref{fig:fitassets}). % src: F:h800_meana_edit_c.tex rows 30-49 (model inside the 95 percent band), 50-69 (model below data), 60; F:h800_mediana_edit_c.tex rows 30-53 (model above data), 54-69 (model below data); 5:303
\citet{Denardi2010} and \citet{lockwood2018incidental} show that a bequest motive and the demand for self-insurance against late-life risks account for saving in old age of the kind the model under-predicts. The under-prediction bears on the saving-subsidy reforms, whose cost and induced income taxes are proportional to the asset base (\autoref{subsec:saving}).

\paragraph{Health information.}
Health is observed by the government: disability insurance pays on the disability state and the earnings condition in \eqref{eq:di}, the Medicare Part B expansion conditions coverage on nonparticipation, and no reform in the menu conditions on health. If health were private information, disability insurance and Medicare would become screening problems in which earnings conditions, asset tests, and waiting periods separate applicants by health, as in \citet{golosov2006designing} and \citet{Low2015}, and the fiscal columns would carry the cost of false claims. The private-information results of the companion paper \citep{ndiaye2026flexible} concern ability, not health.

\paragraph{Time separability.}
Preferences are nonseparable between consumption and leisure within the period and separable across periods (\autoref{subsec:preferences}). Separability over time is a substantive restriction where consumption at the time of exit from the labor force affects the value of consumption afterward: \citet{grochulski2010nonseparable} show that the design of a social security system changes when preferences are nonseparable over time, and \citet{kolsrud2024retirement} use changes in consumption at retirement to evaluate pension design. The welfare valuation of the claiming-schedule reforms, which move the exit age, would differ from \eqref{eq:ce} under such preferences, in a direction that requires a specification.

\paragraph{Income risk.}
Ability follows the AR(1) process of \eqref{eq:wage} with persistence 0.981 and innovation variance 0.016. % src: 11:178; v6:12_app_quant_resubmission_v6.tex:114
Processes with non-Gaussian innovations, heterogeneous profiles, or age-varying variances \citep{storesletten2004consumption,heathcote2010unequal,guvenen2013data} are outside the model. The dispersion of lifetime earnings determines the distribution of average indexed earnings across the brackets of \eqref{eq:pia}, and therefore the incidence of the benefit-formula reforms of \autoref{sec:progressivity}; the persistence of shocks governs the precautionary motive behind the saving and participation responses of every family.

\paragraph{Reentry and irreversibility.}
Exit from the labor force is reversible at the estimated reentry cost of 3,312 hours, at which the model generates no reentry; claiming is irreversible and mandatory by age 70. % src: T:estimated_parameters.tex:16-17; 5:639; F:h800_reentry_edit.tex (model rows all 0); 12:97
Every participation response in the reform tables therefore operates through the timing of exit: an individual who leaves under a steeper early-claiming penalty or a higher normal retirement age does not return, and the 62--64 and 65--69 columns are changes in the stock of individuals who have not yet exited. In the data, reentry rates stay below 5 percent at every age (\autoref{fig:fitassets}); a model that generated them would let part of the response occur through returns to work after exit. % src: F:h800_reentry_edit.tex (data maximum 4.47 at age 63); 11:285
% TODO-AUTHORS: Open question 12 -- confirm that the reform code uses the estimated model (reversible exit with phi_RE, no reentry in equilibrium, irreversible claiming, claiming by 70); if exit is imposed irreversible in the reform code, the first sentence of this paragraph changes.

\subsection{Preference specification and what was not varied}\label{subsec:notvaried}

Appendix~\ref{app:lnc} reports a re-estimation with preferences that are logarithmic in consumption, separable in hours with a Frisch elasticity fixed at 0.5 following \citet{chetty2012bounds}, and with participation and reentry costs in units of utility. % src: 12:710; 12:719
The estimated discount factor is 0.97 and the participation costs are 0.53 for healthy and 0.59 for unhealthy workers (\autoref{tab:estlnc}); as in the baseline, the cost of working is higher for unhealthy than for healthy workers. % src: T:estimated_parameters_lnc.tex:15,20-21; 12:728
The specification matches assets in old age more closely than the baseline, 3.542 against 3.218 in hundreds of thousands of 2016 dollars at age 60, and participation less closely, 96.86 against 83.94 percent for healthy individuals at age 50 (\autoref{fig:fitlnc}). % src: F:h800_lnc_meana.tex row 60; F:h800_l0_lnc_edit.tex row 50; 12:740-746
The reform tables were not re-run under this specification, so the paper does not report how the welfare and fiscal columns depend on the curvature of utility or on the form of the participation cost.

The remaining inputs were not varied. Every reform table is computed at one parameter vector, the point estimates of \autoref{tab:est}, and reports point values, without propagating the sampling uncertainty of that table. The net present values are discounted at one rate, the net return of 3.2 percent; the welfare measure is taken at one age, 30, over one horizon, to age 85; and the dollar conversion uses one consumption window per family (\autoref{sec:welfare}). % src: 5:390; 5:23; 12:632; 12:639
The sensitivities the tables contain are two: the choice of fiscal denominator, net or direct, which gives the lower bounds of \autoref{sec:mvpf} and coincides in sign with the net reading in every row; and the pair of bounds for the saving subsidies, which differ in whether the induced asset income is taxed as it accrues (\autoref{subsec:saving}). Neither is a sensitivity to preferences.
